\chapter{Introduction} \label{chap:introduction}

Kloud Motors is a cloud computing project that implements an online car dealer platform over a large vehicle-listing dataset. The goal of the project is not only to expose marketplace data through an HTTP API, but also to design and operate the platform as a \textbf{cloud-native system} composed of independent services, databases, deployment scripts, and infrastructure configuration. The application therefore combines a realistic business domain, \textit{used and new car sales}, with the technical concerns expected from distributed cloud applications: service decomposition, database isolation, public API design, real-time communication, observability, containerization, Kubernetes deployment, and automated infrastructure support.

\section{Motivation}

Vehicle marketplaces are naturally data-intensive systems. A user looking for a car rarely needs only a static catalogue entry: they compare similar listings, filter by price or mileage, inspect the seller, check whether prices vary by region, save interesting options, and sometimes interact directly with the seller before deciding. Sellers, on the other hand, need ways to expose listings, manage their lifecycle, and potentially use alternative sale formats such as auctions. These requirements make the domain a suitable case study for a cloud platform because it contains both read-heavy analytical workloads and interactive user workflows.

The project was motivated by this combination of marketplace functionality and cloud engineering challenges. A monolithic implementation would be simpler, but it would hide many of the concerns that become important when an application must \textbf{scale, evolve, and tolerate failures}. Kloud Motors was therefore designed around a \textbf{REST and WebSocket gateway} backed by internal gRPC microservices. This makes it possible to separate responsibilities such as listing search, market analysis, user accounts and authentication, seller profiles, chat, and auctions while preserving a single public entry point for clients.

\section{Dataset Characterization}

The platform is based on the \textit{Large Car Dataset} published on Kaggle, a dataset containing used and new car marketplace listings. The dataset has an approximate size of \texttt{5.3 GB} and was published in 2020, with updates in 2021. Its domain fits the project because each record represents the kind of information expected in a car marketplace: vehicle make and model, year, price, mileage, fuel type, location, seller-related information, technical specifications, and listing state.

In the project workflow, the dataset is prepared locally through the scripts in \texttt{scripts/local/} and then loaded into the PostgreSQL databases used by the services. During preparation, the original columns are normalised into fields such as \texttt{VIN}, price, mileage, new/sold state, brand, model, model year, body class, fuel type, transmission, drive type, engine data, colour, seller identifier, and generated location fields. The listing database acts as the \textbf{main source of truth} for catalogue browsing, search, listing details, comparison, market price calculations, and geographical market insights. This structure allows the same base data to support different user facing operations: simple retrieval for listing pages, filtered retrieval for search, aggregation for market analysis, and grouping by city, state, country, or district fields for regional insights.

\section{Business Capabilities}

Kloud Motors exposes a set of business capabilities that together model an online car marketplace. The \textbf{catalogue} capability allows users to retrieve detailed vehicle listings and compare multiple cars side by side. The \textbf{search} capability supports filters such as make, model, year, price range, mileage, fuel type, pagination, and sold-state inclusion, enabling users to narrow a large dataset into relevant results.

The platform also includes analytical capabilities. \textbf{Market price analysis} computes average price information for a selected brand and model, optionally restricted by year ranges. \textbf{Geographical market insights} compare prices and availability across supported location fields, giving users a view of how the same vehicle behaves in different markets. These capabilities turn the dataset from a simple list of cars into decision-support information for buyers and sellers.

Beyond catalogue and analysis features, the application supports marketplace interaction. Seller profiling exposes information about the actor behind a listing. User registration, login, and token refresh are handled by the user service through Firebase authentication, while the gateway validates Firebase ID tokens before forwarding protected requests. Registered users can manage favourite listings. Buyer-seller communication is implemented through chat sessions associated with listings, including WebSocket-based real-time messaging. Chat indexing is stored in PostgreSQL, message persistence is backed by Firestore, and Google Pub/Sub can distribute real-time events across replicas. The auction module adds another transactional workflow by allowing authenticated sellers to create auctions and authenticated buyers to place bids while clients receive live auction updates, also with Pub/Sub support for distributed WebSocket updates.

\section{Use Cases}

The main use cases follow directly from these capabilities. A visitor can search the catalogue using filters and inspect detailed listings before comparing selected cars. A buyer interested in fair pricing can request average market prices for a specific brand and model, while another user can compare prices by location to understand regional differences. When a listing is promising, the visitor can inspect the seller profile, register or log in through the \texttt{/api/user} routes, save the listing as a favourite through the protected \texttt{/api/users/me/favorites} routes, and open a chat with the seller.

The platform also supports seller-oriented and real-time scenarios. Authenticated users can create, update, or delete listings through the gateway. Sellers can create auctions for specific listings, buyers can place bids, and authenticated clients can receive auction updates over WebSockets. These use cases exercise both the data-oriented and interactive parts of the system, making Kloud Motors a representative cloud application rather than a purely static API over a dataset.

The remainder of this report presents the API contract, functional and non-functional requirements, service architecture, implementation details, deployment model, pipeline, validation strategy, future improvements, and cost analysis of the Kloud Motors platform.
